% ============================================================================
% Fichier  : partie1/P01_C01_philosophie_fondements.tex
% Titre    : Philosophie et Fondements de l'Investigation Numérique
% Auteur   : MINKA MI NGUIDJOÏ Thierry Emmanuel
% Version  : 1.0 -- Juin 2026
% Statut   : RÉDIGÉ
%
% Sources fusionnées :
%   - 0_philosophie_fondements.tex (version précédente)
%   - chap2.pdf (archéologie foucaldienne des régimes de vérité)
%   - minka2025cro (Trilemme CRO, ePrint 2025/1348)
%   - minka2025zknr (paradoxe de l'authenticité invisible, ePrint 2025/1138)
% ============================================================================

\chapter{Philosophie et Fondements de l'Investigation Numérique}
\label{chap:phi-fond}

\epigraph{%
    « L'archéologie ne cherche pas à retrouver la continuité ininterrompue~;
    elle établit ce qu'il nous est possible de connaître. »%
}{--- Michel Foucault, \textit{L'Archéologie du savoir}, 1969}

\niveauChapitre{Fondamental}{%
    Aucun prérequis. Ce chapitre articule philosophie, épistémologie et
    mathématiques à un niveau accessible. Les équations servent à préciser
    des intuitions, non à les remplacer.%
}

\begin{boxobjectifs}
\begin{itemize}
    \item Caractériser l'investigation numérique comme pratique philosophique
          à part entière, distincte d'une simple maîtrise d'outils.
    \item Comprendre la notion foucaldienne de \termecle{régime de vérité
          numérique} et l'appliquer à la périodisation de la discipline.
    \item Maîtriser les fondements mathématiques mobilisés en forensique~:
          entropie de Shannon, théorie des graphes, sensibilité aux
          conditions initiales.
    \item Articuler le \termecle{paradoxe de l'authenticité invisible}
          avec le Trilemme CRO comme grille de lecture unifiée.
    \item Formuler les obligations éthiques de l'investigateur numérique
          à partir de la notion de \termecle{praxis de liberté}.
\end{itemize}
\end{boxobjectifs}

\bigskip

% ============================================================================
\section{L'Investigation Numérique comme Pratique Discursive}
% ============================================================================

\subsection{Au-delà de la technique : une discipline philosophique}

L'investigation numérique est souvent présentée comme une affaire d'outils~:
Autopsy pour analyser les disques, Volatility pour la mémoire, Wireshark pour
le réseau. Cette réduction est non seulement inexacte -- elle est dangereuse.
Un investigateur réduit à ses outils est un artisan. Ce cours forme des
praticiens capables de produire des conclusions défendables devant n'importe
quel tribunal, dans n'importe quel contexte, avec les outils disponibles au
moment de l'investigation. Cela exige une \textbf{théorie}, pas seulement
une boîte à outils.

\medskip

La théorie commence ici, avec une question que Foucault aurait reconnue~:
\textit{qu'est-ce qui fait qu'un énoncé numérique compte comme preuve à
une époque donnée ?} Cette question n'est pas rhétorique. La réponse a changé
quatre fois depuis 1970, et elle changera encore. L'investigateur qui ne
comprend pas ces changements travaille dans l'obscurité~; celui qui les
comprend peut anticiper les défis de l'audience avant même de toucher
le premier disque dur.

\begin{boxdef}
\textbf{Investigation numérique}\index{investigation numérique} (\textit{Digital forensics}) ---
Discipline scientifique appliquée ayant pour objet l'identification,
la préservation, l'analyse et la présentation de preuves numériques en
respectant des méthodes juridiquement admissibles. Elle mobilise
simultanément l'informatique, le droit, les mathématiques et la philosophie
de la connaissance.
\end{boxdef}

\subsection{La notion foucaldienne de régime de vérité numérique}

Notre cadre d'analyse s'inspire de \textit{L'Archéologie du savoir} de
Foucault \cite{foucault1969} pour analyser l'investigation numérique comme
\termecle{pratique discursive} productrice de vérité. Chaque période
historique correspond à un \textbf{régime de vérité numérique} différent~:
un système qui détermine a priori ce qui peut compter comme preuve légitime,
quelles méthodes d'investigation sont autorisées, quelles institutions sont
habilitées à certifier la vérité numérique, et quel statut épistémique est
accordé à l'investigateur.

\begin{boxdef}
\textbf{Régime de vérité numérique}\index{régime de vérité numérique} --- Système
historiquement situé qui détermine, de manière a priori~: les types de
traces considérées comme preuves légitimes~; les techniques d'investigation
autorisées~; les institutions habilitées à certifier~; les conditions
d'acceptabilité juridique et sociale~; et le statut épistémique de
l'investigateur comme sujet de savoir.
\end{boxdef}

Cette notion est plus qu'un exercice académique. Elle explique pourquoi
certaines preuves irréfutables sur le plan technique restent irrecevables
sur le plan juridique~: elles appartiennent à un régime de vérité que le
droit n'a pas encore assimilé. Elle explique aussi pourquoi un investigateur
formé uniquement sur les outils actuels sera dépassé dans dix ans, tandis
que celui qui comprend la logique des régimes pourra s'adapter.

\subsection{Périodisation mathématique des régimes}

La périodisation des régimes de vérité repose sur quatre axes d'analyse
dont la dominance relative définit chaque époque. Pour toute période
historique $t$, le régime est caractérisé par un \textbf{vecteur de
dominance} \cite{minka2025cro}~:

\begin{equation}
\vec{R}_t = \bigl(\alpha_T,\; \alpha_J,\; \alpha_S,\; \alpha_P\bigr)
\quad \text{avec} \quad \sum_{i} \alpha_i = 1
\label{eq:vecteur-regime}
\end{equation}

où $\alpha_T$ représente la dominance technologique, $\alpha_J$ la
dominance juridique et normative, $\alpha_S$ la dominance sociale et
culturelle, et $\alpha_P$ la dominance des pratiques professionnelles.
La transition d'un régime $\vec{R}_t$ à son successeur $\vec{R}_{t+1}$
obéit à une équation de transition non linéaire~:

\begin{equation}
\vec{R}_{t+1} = F\bigl(\vec{R}_t,\; \Delta\mathrm{Tech}_t,\;
\Delta\mathrm{Legal}_t,\; \mathcal{I}_t\bigr)
\label{eq:transition}
\end{equation}

où $\mathcal{I}_t$ représente l'ensemble des incidents critiques de la
période $t$ -- les affaires qui ont contraint le régime à se reconfigurer.

\bigskip

Le tableau suivant synthétise l'évolution des régimes depuis l'émergence
de la discipline \cite{chap2pdf}~:

\begin{tableaucours}{p{2.0cm}p{2.8cm}p{3.2cm}p{3.4cm}}{%
    \tableauHeader{Période} &
    \tableauHeader{Régime dominant} &
    \tableauHeader{Preuve paradigmatique} &
    \tableauHeader{Autorité épistémique}
}
1970--1990 & Technique $(\alpha_T \approx 0.7)$ & Log système, adresse IP &
    Expert technique isolé \\
1990--2000 & Juridique $(\alpha_J \approx 0.5)$ & Chaîne de custody & Tribunal \\
2000--2010 & Standardisé $(\alpha_P \approx 0.5)$ & Conformité ISO/IEC & Organisme
    de certification \\
2010--2020 & Computationnel $(\alpha_T \approx 0.6)$ & Analyse big data &
    Algorithme \\
2020--\ldots & Quantique / ZK & Preuve à divulgation nulle &
    \sigle{IA} / Protocole \\
\end{tableaucours}
\vspace{-0.8em}
\begin{center}
{\small\itshape Tableau~1.1 --- Évolution des régimes de vérité numérique
(d'après \cite{chap2pdf} et \cite{minka2025cro})}
\end{center}

% ============================================================================
\section{Ontologie de la Trace Numérique}
% ============================================================================

\subsection{La trace comme phénomène existentiel}

Dans la philosophie de Derrida, la \textit{trace} est l'inscription d'une
absence~: quelque chose était là, a agi, et a laissé une marque. La
\termecle{trace numérique}\index{trace numérique} est une instanciation
particulièrement pure de cette notion~: elle est produite sans intention
consciente de laisser une marque, elle persiste au-delà de l'action qui
l'a générée, et elle peut être interprétée comme révélatrice d'intention
\emph{a posteriori}.

\medskip

La trace numérique n'est pas une simple donnée~; elle est une
\textbf{manifestation d'existence}. Chaque fichier créé, chaque connexion
établie, chaque frappe de clavier enregistrée dans un journal est la
matérialisation numérique d'un acte humain. L'investigateur qui lit ces
traces ne lit pas de la technique~: il lit du comportement, de l'intention,
de la relation. Cette lecture est herméneutique avant d'être technique.

\begin{boiteRemarque}{Herméneutique des données}
L'interprétation des traces numériques requiert ce que Gadamer appelle
un \textit{cercle herméneutique}~: on ne comprend une trace qu'à la lumière
du contexte global, et on ne comprend le contexte global qu'à la lumière
des traces individuelles. En pratique~: un fichier \texttt{.docx} nommé
\texttt{faux\_bon\_de\_commande.docx} n'a de signification forensique que
replacé dans la chronologie des virements, l'historique de navigation et
les échanges d'emails. Chaque artefact éclaire les autres.
\end{boiteRemarque}

\subsection{Temporalité numérique multidimensionnelle}

Le temps numérique n'est pas linéaire. Un fichier possède jusqu'à quatre
horodatages distincts dans le système de fichiers NTFS~:
\textbf{M}odified, \textbf{A}ccessed, \textbf{C}hanged, \textbf{B}orn
(horodatage de création). Ces quatre dimensions peuvent diverger de
façon significative~: un fichier créé en 2020, modifié en 2024, mais
accédé pour la dernière fois en 2022 raconte une histoire non linéaire.
L'investigateur doit maîtriser cette \textbf{temporalité
plurielle}\index{temporalité numérique} pour ne pas
commettre d'erreurs d'interprétation.

\begin{tableaucours}{p{2.2cm}p{3.8cm}p{5.2cm}}{%
    \tableauHeader{Horodatage} &
    \tableauHeader{Signification} &
    \tableauHeader{Remarque forensique}
}
\texttt{Modified} & Dernière modification du contenu & Le plus couramment
    analysé~; peut être falsifié (\textit{timestomping}) \\
\texttt{Accessed} & Dernier accès en lecture & Peu fiable~: Windows désactive
    souvent la mise à jour pour des raisons de performance \\
\texttt{Changed} & Dernière modification des métadonnées & Révèle les
    opérations sur l'entrée MFT~; difficile à falsifier \\
\texttt{Born} & Date de création sur ce volume & Clé pour détecter les
    fichiers copiés (Born > Modified = copié depuis un autre volume) \\
\end{tableaucours}
\vspace{-0.8em}
\begin{center}
{\small\itshape Tableau~1.2 --- Les quatre horodatages NTFS (timestamps MAC-B)
et leur interprétation forensique}
\end{center}

La comparaison entre la preuve traditionnelle et la preuve numérique révèle
l'ampleur de cette transformation ontologique~:

\begin{tableaucours}{p{5.5cm}p{5.5cm}}{%
    \tableauHeader{Preuve traditionnelle} &
    \tableauHeader{Preuve numérique}
}
Matérialité tangible, physique & Immatérialité des bits \\
Stabilité physique dans le temps & Volatilité et mutabilité \\
Authenticité portée par l'objet & Authenticité par la chaîne de confiance \\
Temporalité linéaire, univoque & Temporalité multidimensionnelle (MAC-B) \\
Falsification visible à l'expert & Falsification indétectable sans hachage \\
Localité géographique claire & Dispersion cloud multi-juridictionnelle \\
\end{tableaucours}
\vspace{-0.8em}
\begin{center}
{\small\itshape Tableau~1.3 --- Transition ontologique de la preuve~:
du physique au numérique}
\end{center}

% ============================================================================
\section{Épistémologie de la Preuve Numérique}
% ============================================================================

\subsection{La crise de la vérité numérique}

L'ère numérique engendre une \textbf{crise de la vérité} sans précédent.
Trois phénomènes convergents l'alimentent~:

\begin{itemize}
    \item \textbf{Manipulation algorithmique~:} les \textit{deepfakes}, les
          documents générés par intelligence artificielle, les voix synthétiques
          brouillent la frontière entre l'authentique et le fabriqué avec une
          précision qui dépasse la capacité de détection humaine non assistée.

    \item \textbf{Érosion de l'autorité épistémique~:} la multiplication des
          sources d'information, la désintermédiation des médias et la
          fragmentation des espaces discursifs ont détruit le consensus sur
          ce qui constitue une source fiable.

    \item \textbf{Fragmentation du réel~:} les mêmes événements peuvent être
          documentés par des enregistrements contradictoires, tous authentiques,
          correspondant à des points de vue partiels sur une réalité complexe.
\end{itemize}

L'investigateur numérique devient dans ce contexte un
\textbf{archiviste du réel}\index{archiviste du réel}~: celui dont le rôle
est de préserver l'intégrité de la mémoire collective contre l'effacement,
la falsification et la réécriture. Ce rôle dépasse largement le cadre
judiciaire~; il touche à la démocratie, à la justice et à l'histoire.

\subsection{La théorie de Kuhn appliquée à la forensique}

Thomas Kuhn, dans \textit{La Structure des révolutions scientifiques}
\cite{kuhn1962}, a montré que les sciences ne progressent pas par accumulation
continue mais par ruptures paradigmatiques. La forensique numérique illustre
parfaitement ce modèle~: chaque décennie depuis 1970 a correspondu à un
changement de paradigme, pas à une simple amélioration.

\medskip

Le paradigme \textbf{technique} (1970--1990) posait que la vérité numérique
était une affaire d'experts isolés lisant des logs. Le paradigme
\textbf{juridique} (1990--2000) a introduit la chaîne de custody et fait
du tribunal l'arbitre final. Le paradigme \textbf{standardisé} (2000--2010)
a délégué cette autorité aux organismes de certification (\sigle{ISO/IEC}).
Le paradigme \textbf{computationnel} (2010--2020) a transféré l'autorité
aux algorithmes d'analyse massive. Le paradigme \textbf{quantique/ZK}
(2020--\ldots) en cours de formation pose que la vérité numérique peut et
doit être prouvée sans divulgation de contenu.

\medskip

Chaque transition a rendu caduques des pans entiers de la pratique
existante. Un investigateur formé uniquement dans le paradigme standardisé
-- maîtrisant parfaitement les procédures ISO~27037 -- n'est pas
naturellement préparé à travailler avec des preuves à divulgation nulle.
C'est pourquoi ce cours enseigne les \emph{fondements}, pas les procédures~:
les fondements survivent aux changements de paradigme.

% ============================================================================
\section{Fondements Mathématiques de la Discipline}
% ============================================================================

Les mathématiques de l'investigation numérique ne sont pas décoratives~:
elles fournissent des outils opérationnels utilisés quotidiennement.
Trois théories fondent la discipline.

\subsection{Théorie de l'information : l'entropie de Shannon}
\label{subsec:shannon}

La théorie de l'information de Shannon \cite{shannon1948} quantifie
l'incertitude contenue dans un système d'information. Son concept central,
l'\termecle{entropie informationnelle}\index{entropie de Shannon}, est
défini par~:

\begin{equation}
H(X) = -\sum_{i=1}^{n} P(x_i) \log_2 P(x_i)
\label{eq:shannon}
\end{equation}

où $X$ est une variable aléatoire prenant les valeurs $x_1, \ldots, x_n$
avec les probabilités $P(x_1), \ldots, P(x_n)$, et $H(X)$ est mesurée
en bits. Cette équation a trois applications forensiques directes~:

\begin{enumerate}
    \item \textbf{Détection de données chiffrées ou compressées~:}
          un fichier chiffré ou compressé présente une entropie proche de~8
          bits/octet (distribution quasi uniforme des octets). Un fichier texte
          ordinaire présente une entropie de~3 à~5 bits/octet. Un examinateur
          peut ainsi identifier rapidement les zones d'un disque contenant
          des données potentiellement chiffrées.

    \item \textbf{Détection de stéganographie~:} l'insertion de données cachées
          dans un fichier image ou audio modifie la distribution statistique
          des pixels ou des échantillons, entraînant une divergence entropique
          mesurable par rapport à un fichier non modifié de même type.

    \item \textbf{Analyse de malware~:} les sections de code d'un exécutable
          malveillant chiffré (shellcode chiffré, packer) présentent une
          entropie anormalement haute, ce qui permet de les identifier avant
          même d'exécuter le binaire.
\end{enumerate}

\begin{boxartefact}{Entropie d'un fichier -- indicateur forensique}
\textbf{Interprétation pratique de l'entropie~:}\\[0.5ex]
\begin{tabular}{p{3cm}p{7.5cm}}
$H \in [0, 1]$ & Fichier quasi vide ou très répétitif (logs compressés, fichiers
    sparse) \\
$H \in [3, 5]$ & Texte naturel, code source, XML/JSON non compressé \\
$H \in [6, 7]$ & Données binaires compilées, exécutables standard \\
$H \in [7.5, 8]$ & Chiffrement, compression, stéganographie -- investigation
    approfondie recommandée \\
\end{tabular}
\end{boxartefact}

\subsection{Théorie des graphes : modélisation des relations}

L'analyse relationnelle repose sur la \termecle{théorie des graphes},
qui modélise les entités et leurs interactions~:

\begin{equation}
G = (V, E) \quad \text{avec} \quad E \subseteq V \times V
\label{eq:graphe}
\end{equation}

où $V$ est l'ensemble des \textit{sommets} (entités~: personnes, machines,
comptes, adresses IP) et $E$ l'ensemble des \textit{arêtes} (relations~:
communications, transferts, accès). Un graphe pondéré $G = (V, E, w)$ où
$w : E \to \mathbb{R}$ associe un poids à chaque relation (volume de
communications, fréquence, montants transférés) est particulièrement
utile en investigation financière.

\medskip

Les applications forensiques sont nombreuses~:

\begin{itemize}
    \item \textbf{Analyse de réseau de communication~:} les échanges d'emails
          ou de messages forment un graphe dont les algorithmes de centralité
          (PageRank, betweenness centrality) permettent d'identifier les
          acteurs clés d'un réseau criminel.
    \item \textbf{Corrélation d'incidents~:} les logs de plusieurs systèmes
          peuvent être modélisés comme un graphe temporel pour reconstituer
          un chemin d'attaque.
    \item \textbf{Forensique blockchain~:} les transactions
          cryptomonnaies forment un graphe de flux dont l'analyse permet
          de remonter des fonds jusqu'à des entités identifiables.
\end{itemize}

\subsection{Sensibilité aux conditions initiales}

L'investigation numérique opère dans des systèmes complexes où de minuscules
altérations peuvent avoir des conséquences considérables. Cela est formalisé
par la théorie du chaos~: pour un système dynamique, la distance entre
deux trajectoires initialement proches croît exponentiellement~:

\begin{equation}
\delta(t) \approx \delta(0) \cdot e^{\lambda t}
\label{eq:lyapunov}
\end{equation}

où $\lambda$ est l'exposant de Lyapunov caractérisant la sensibilité du
système. En forensique, ce résultat justifie mathématiquement l'obsession
pour la préservation de l'intégrité~: une modification infinitésimale de
l'état initial (allumer le PC, monter une partition) peut produire
une divergence massive dans l'état final (modification de centaines
d'horodatages, création de fichiers temporaires, mise à jour du registre).
La sensibilité aux conditions initiales est la justification théorique
profonde du principe d'\textbf{intégrité}.

% ============================================================================
\section{Le Paradoxe de l'Authenticité Invisible}
% ============================================================================

\subsection{Énoncé du paradoxe}

Le \termecle{paradoxe de l'authenticité invisible}\index{paradoxe de
l'authenticité invisible} -- théorisé dans \cite{minka2025zknr} -- capture
la tension fondamentale entre authentification et confidentialité~:

\begin{quotation}
\textit{Plus une preuve numérique est authentique et vérifiable, plus elle
tend à révéler des informations sur son contenu et son origine, compromettant
ainsi la confidentialité. Inversement, plus une preuve préserve la
confidentialité, plus son authenticité devient difficile à établir de
manière certaine.}
\end{quotation}

Cette tension peut s'exprimer comme une relation d'incertitude analogue
au principe d'Heisenberg~:

\begin{equation}
\Delta\mathcal{A} \cdot \Delta\mathcal{C} \;\geq\; \hbar_{\mathrm{num}}
\label{eq:authenticite-invisible}
\end{equation}

où $\Delta\mathcal{A}$ est l'incertitude sur l'authenticité (fiabilité de
la preuve), $\Delta\mathcal{C}$ est l'incertitude sur la confidentialité
(protection de l'information), et $\hbar_{\mathrm{num}}$ est une constante
numérique fondamentale caractérisant le système considéré.

\begin{boiteRemarque}{Statut de la formule~\eqref{eq:authenticite-invisible}}
Cette relation d'incertitude est une \emph{analogie formalisée}, pas une
dérivation de la mécanique quantique. Sa valeur est heuristique~: elle
nomme une tension réelle et irréductible, lui donne une forme mathématique
précise, et ouvre la voie à des mesures opérationnelles. La constante
$\hbar_{\mathrm{num}}$ dépend du système cryptographique et du contexte
juridique~; sa quantification précise est l'objet de travaux en cours
au LIMSI.
\end{boiteRemarque}

\subsection{Implications épistémologiques}

Le paradoxe soulève trois questions profondes que l'investigateur doit
intégrer dans sa pratique~:

\begin{enumerate}
    \item \textbf{Peut-on savoir qu'une preuve est authentique sans
          en connaître le contenu~?} Les preuves à divulgation nulle
          (protocoles ZK) répondent oui -- à condition d'accepter que
          la preuve change de nature~: elle n'est plus un objet à examiner
          mais un processus à vérifier.

    \item \textbf{Comment fonder la confiance dans ce qui reste invisible~?}
          La réponse est institutionnelle autant que technique~: la confiance
          repose sur la robustesse mathématique du protocole ET sur la
          légitimité de l'institution qui le met en œuvre.

    \item \textbf{L'authenticité peut-elle exister sous forme chiffrée~?}
          Oui~: un hash cryptographique prouve l'intégrité d'un fichier
          sans en révéler le contenu. Mais l'opposabilité juridique de
          cette preuve reste dépendante de la juridiction.
\end{enumerate}

\subsection{Résolution partielle : les protocoles ZK-NR}

Les protocoles \termecle{Zero-Knowledge Non-Repudiation} (ZK-NR)
\cite{minka2025zknr} apportent une résolution pratique en permettant~:

\[
\text{Vérification} \quad \text{sans} \quad \text{Divulgation}
\]
\[
\text{Confiance} \quad \text{sans} \quad \text{Transparence totale}
\]
\[
\text{Preuve} \quad \text{sans} \quad \text{Révélation du contenu}
\]

Cette résolution n'est pas magique~: elle ne supprime pas la tension
du paradoxe, elle en gère les effets de façon contrôlée. La Partie~VI
de ce cours (Chapitres~25--31) est entièrement consacrée à ces protocoles.
Le paradoxe de l'authenticité invisible sera alors revisité dans toute
sa profondeur mathématique.

% ============================================================================
\section{Le Trilemme CRO comme Grille de Lecture Unifiée}
% ============================================================================

\subsection{Genèse et formalisation}

Les tensions identifiées dans les sections précédentes -- entre
authentification et confidentialité, entre opposabilité et protection de
la vie privée, entre rigueur technique et admissibilité juridique -- ne
sont pas des accidents. Elles constituent la structure profonde de toute
investigation numérique, formalisée par le \termecle{Trilemme CRO}
\cite{minka2025cro}.

\begin{boxCRO}
\textbf{Le Trilemme CRO --- Définition formelle}

\medskip
Soit $\Pi$ un protocole ou une méthode d'investigation numérique. On définit~:
\begin{itemize}
    \item $\mathcal{C}(\Pi) \in [0,1]$~: indice de confidentialité ---
          mesure dans quelle proportion $\Pi$ protège les informations
          sensibles des parties impliquées.
    \item $\mathcal{R}(\Pi) \in [0,1]$~: indice de fiabilité ---
          mesure dans quelle proportion $\Pi$ garantit l'intégrité et
          la reproductibilité des preuves.
    \item $\mathcal{O}(\Pi) \in [0,1]$~: indice d'opposabilité juridique ---
          mesure dans quelle proportion $\Pi$ produit des preuves
          recevables et défendables devant un tribunal.
\end{itemize}

\textbf{Énoncé du trilemme \cite{minka2025cro}~:}
\begin{equation}
\Gamma_{CRO}(\Pi) \;=\; \max\bigl\{\mathcal{C}(\Pi),\;
\mathcal{R}(\Pi),\; \mathcal{O}(\Pi)\bigr\}
\;\geq\; 0.4 + \mathrm{negl}(\lambda)
\label{eq:CRO}
\end{equation}

où $\lambda$ est le paramètre de sécurité du système cryptographique sous-jacent
et $\mathrm{negl}(\lambda)$ une fonction négligeable.

\medskip
\textit{Interprétation~: aucun protocole ne peut simultanément maximiser
les trois indices. Au moins l'une des trois dimensions sera contrainte à
demeurer sous son maximum théorique. La compétence de l'investigateur
se mesure précisément à sa capacité à choisir et justifier
le compromis optimal pour le contexte.}
\end{boxCRO}

\subsection{Le Trilemme comme boussole de la pratique}

À chaque étape d'une investigation, l'analyste du Trilemme CRO pose
trois questions~:

\begin{enumerate}
    \item \textbf{Cette technique augmente-t-elle $\mathcal{C}$, $\mathcal{R}$
          ou $\mathcal{O}$ ?} Si elle augmente $\mathcal{R}$ (ex.~: capture
          RAM complète), quelle dimension sacrifie-t-elle ?
          ($\mathcal{C}$~: accès à des données personnelles non pertinentes).

    \item \textbf{Le compromis est-il justifiable et documenté~?} Un
          compromis non documenté n'est pas défendable. Un compromis documenté
          et justifié est presque toujours défendable.

    \item \textbf{Ce compromis est-il proportionnel à l'enjeu~?} Une
          investigation pour un détournement de 500~000~FCFA ne justifie
          pas la même atteinte à $\mathcal{C}$ qu'une investigation pour
          une affaire de terrorisme.
\end{enumerate}

Le Trilemme CRO n'est pas une contrainte abstraite~: c'est l'outil qui
transforme chaque décision technique en décision éthique consciente. Il
sera appliqué systématiquement à chaque chapitre de ce cours, sous la
forme de l'encadré \textbf{Analyse~CRO}.

% ============================================================================
\section{Éthique et Responsabilité de l'Investigateur}
% ============================================================================

\subsection{L'investigateur comme philosophe-praticien}

L'investigateur numérique moderne endosse ce que Ricoeur appellerait
une \textbf{identité narrative composite}~: il est simultanément~:

\begin{enumerate}
    \item \textbf{Archéologue du digital~:} il exhume et préserve les traces
          numériques comme un archéologue préserverait des vestiges -- avec
          la conscience que toute intervention modifie irréversiblement ce
          qu'on tente de conserver.

    \item \textbf{Épistémologue pratique~:} il évalue en permanence la
          fiabilité de ce qu'il trouve, distingue la corrélation de la
          causalité, et refuse de conclure au-delà de ce que les preuves
          permettent.

    \item \textbf{Éthicien appliqué~:} il navigue des dilemmes moraux réels,
          souvent sans temps de réflexion, et doit avoir internalisé des
          principes suffisamment robustes pour guider ses choix dans l'urgence.
\end{enumerate}

\subsection{L'investigation comme praxis de liberté}

Bien comprise, l'investigation numérique est une \termecle{praxis de
liberté}\index{praxis de liberté}~: elle protège les individus contre
l'arbitraire en documentant le réel~; elle permet la reddition des comptes
dans des sociétés où le pouvoir numérique est concentré~; elle préserve
la mémoire collective contre l'effacement délibéré~; elle équilibre des
rapports de force asymétriques par la transparence.

\medskip

Cette dimension politique et éthique n'est pas extérieure à la technique~:
elle lui donne son sens. C'est pourquoi ce cours insiste autant sur les
conditions juridiques d'admissibilité que sur les commandes forensiques~:
une preuve obtenue illégalement ou de façon disproportionnée n'est pas
seulement irrecevable -- elle est une atteinte aux droits fondamentaux.

\begin{boxdef}
\textbf{Charte de l'investigateur numérique}\index{charte de l'investigateur}
--- Ensemble des obligations éthiques fondamentales~:

\begin{enumerate}
    \item Je préserverai l'intégrité de la preuve avant toute autre
          considération, y compris celle d'efficacité.
    \item Je respecterai la dignité numérique des personnes impliquées,
          qu'elles soient suspectes, victimes ou témoins.
    \item Je reconnaîtrai les limites de ma connaissance et les documenterai
          dans mes rapports.
    \item Je travaillerai pour la vérité, pas pour la conviction d'une
          partie -- ma loyauté est envers le tribunal, pas envers
          le commanditaire.
    \item Je me souviendrai que derrière chaque donnée se trouve un humain
          dont la vie peut être affectée par mes conclusions.
\end{enumerate}
\end{boxdef}

\subsection{Vers une éthique post-quantique}

L'avènement de la cryptographie post-quantique et des preuves à divulgation
nulle soulève de nouveaux impératifs éthiques que la pratique actuelle
n'a pas encore intégrés. Trois enjeux majeurs~:

\begin{itemize}
    \item \textbf{La résistance à l'obsolescence~:} les preuves numériques
          collectées aujourd'hui avec des algorithmes cryptographiques
          classiques (RSA, ECDSA) seront potentiellement cassables dans dix
          à quinze ans par un ordinateur quantique suffisant. L'investigateur
          a-t-il une obligation de collecter dès maintenant avec des
          algorithmes post-quantiques~? La réponse est oui, pour les
          affaires à haute valeur ou longue durée.

    \item \textbf{Le droit à l'oubli dans un monde de mémoire parfaite~:}
          les capacités croissantes de stockage et de traitement permettent
          de conserver et d'analyser des traces numériques sur des décennies.
          L'investigateur doit intégrer le principe de proportionnalité
          temporelle~: la durée de conservation des preuves collectées doit
          être proportionnelle à l'enjeu judiciaire.

    \item \textbf{La responsabilité des algorithmes~:} quand un outil d'IA
          identifie un suspect ou classe un artefact, qui est responsable
          de cette conclusion~? L'investigateur ne peut pas déléguer
          sa responsabilité épistémique à un algorithme opaque.
\end{itemize}

\separateur

\begin{boxresume}
\textbf{Ce qu'il faut retenir de ce chapitre~:}

\begin{itemize}
    \item L'investigation numérique est une \textbf{pratique discursive}
          au sens foucaldien~: chaque époque définit ses propres critères
          de vérité numérique. Comprendre ces régimes permet d'anticiper
          les défis judiciaires à venir.

    \item Le \textbf{vecteur de dominance} $\vec{R}_t = (\alpha_T, \alpha_J,
          \alpha_S, \alpha_P)$ formalise mathématiquement le régime de vérité
          dominant à chaque période.

    \item La \textbf{trace numérique} est une manifestation d'existence~:
          sa lecture est herméneutique avant d'être technique. La temporalité
          numérique est multidimensionnelle (horodatages MAC-B).

    \item Trois outils mathématiques fondent la discipline~:
          \textbf{entropie de Shannon} (détection d'anomalies),
          \textbf{théorie des graphes} (analyse relationnelle),
          \textbf{sensibilité aux conditions initiales} (justification
          du principe d'intégrité).

    \item Le \textbf{paradoxe de l'authenticité invisible}
          ($\Delta\mathcal{A} \cdot \Delta\mathcal{C} \geq \hbar_{\mathrm{num}}$)
          nomme la tension irréductible entre authentification et
          confidentialité. Les protocoles ZK-NR en gèrent les effets,
          sans la supprimer.

    \item Le \textbf{Trilemme CRO} ($\Gamma_{CRO}(\Pi) \geq 0.4 +
          \mathrm{negl}(\lambda)$) est la grille d'analyse qui traverse
          tout ce cours~: toute investigation est un compromis explicite
          entre Confidentialité, Fiabilité et Opposabilité.

    \item L'investigateur est un \textbf{philosophe-praticien}~: sa
          responsabilité est envers la vérité et la justice, pas envers
          la conviction d'une partie. La Charte de l'investigateur
          traduit cette responsabilité en obligations concrètes.
\end{itemize}
\end{boxresume}

\bigskip

\begin{boxexo}[title={\ding{45}\quad Exercice~1.1 --- Diagnostic de régime \niveaubase}]
Identifiez le régime de vérité numérique dominant (1970--1990,
1990--2000, 2000--2010, 2010--2020 ou 2020--\ldots) pour chacun des
scénarios suivants, en justifiant votre réponse par les caractéristiques
du régime (preuve paradigmatique, autorité épistémique, vecteur dominant)~:

\begin{enumerate}[(a)]
    \item Un juge rejette une preuve numérique au motif qu'elle a été
          collectée sans protocole de hachage documenté.
    \item Un investigateur présente au tribunal une analyse de clustering
          automatique de 2~millions d'emails, sans pouvoir expliquer
          précisément comment l'algorithme a identifié les messages suspects.
    \item Une entreprise engage un consultant qui, à partir de logs système,
          reconstitue manuellement l'activité d'un intrus et rédige
          un rapport non structuré mais techniquement exact.
    \item Une unité spécialisée utilise une preuve à divulgation nulle
          pour prouver qu'un document a été modifié, sans révéler son contenu.
\end{enumerate}
\end{boxexo}

\begin{boxexo}[title={\ding{45}\quad Exercice~1.2 --- Application du Trilemme CRO \niveaumoyen}]
Pour chacune des situations suivantes, calculez qualitativement les trois
indices CRO (\textit{faible / moyen / élevé}), identifiez le compromis
réalisé et évaluez s'il est justifiable~:

\begin{enumerate}[(a)]
    \item Un investigateur capture l'intégralité du trafic réseau d'une
          entreprise pendant 72~heures pour analyser une fuite de données.
    \item Un OPJ exige du prestataire cloud qu'il livre les données chiffrées
          d'un suspect sans lui fournir la clé de déchiffrement.
    \item Un analyste publie ses outils d'investigation en open source pour
          garantir la reproductibilité de ses méthodes.
    \item Une juridiction autorise l'utilisation d'une analyse IA pour
          classer les messages comme ``suspects'' mais interdit à l'avocat
          de la défense d'accéder à l'algorithme utilisé.
\end{enumerate}
\end{boxexo}

\begin{boxexo}[title={\ding{45}\quad Exercice~1.3 --- Réflexion sur l'entropie \niveauexpert}]
Vous analysez une image forensique d'un disque dur de 500~Go. L'outil
d'analyse d'entropie détecte une zone de 47~Go présentant une entropie
uniforme de~7.98~bits/octet. Le reste du disque présente des entropies
variables entre 3.2 et 6.8~bits/octet.

\begin{enumerate}[(a)]
    \item Quelle est votre hypothèse principale sur la nature de cette zone~?
    \item Proposez deux hypothèses alternatives (sans supposer d'activité
          malveillante) qui pourraient expliquer cette observation.
    \item Quelles analyses complémentaires effectueriez-vous pour discriminer
          entre ces hypothèses~? Précisez les outils et les commandes.
    \item Comment documenteriez-vous ces analyses dans votre rapport pour
          satisfaire simultanément $\mathcal{R}$ et $\mathcal{O}$~?
\end{enumerate}
\end{boxexo}

\bigskip

\begin{boxinfo}
\textbf{Ce qui vient.} Le Chapitre~2 reprendra le cadre foucaldien posé
ici pour décrire, affaire par affaire, comment chaque régime de vérité
s'est constitué historiquement. La périodisation mathématique
$\vec{R}_t$ servira de fil conducteur. Les fondements mathématiques
(entropie, graphes) seront mis en pratique dans la Partie~IV. Le Trilemme
CRO et le paradoxe de l'authenticité invisible seront formalisés
rigoureusement dans la Partie~VI.
\end{boxinfo}
